\section{字符级GPT构建与预训练实验}

本节将前述机制落实为网络配置、训练样本与优化程序，并通过验证损失、生成示例及受控迁移实验评价预训练效果。

\Needspace{8\baselineskip}
\subsection{模型配置与计算结构}

生成主干使用PyTorch基本神经网络模块构建，并从随机初始化训练。实验设备为RTX 5060 Laptop 8GB。模型设置6层、隐藏维256、4个注意力头和1024维前馈层，包含8,154,624个参数。采用该规模的目的是在单卡预算内保留完整自回归Transformer计算流程。表~\ref{tab:models}汇总基础模型与两个任务模型的配置。

\begin{table}[htbp]
 \centering\small
 \caption{基础模型与两个任务模型的主要配置}\label{tab:models}
 \begin{tabularx}{\linewidth}{Yccc}
 \toprule 配置项&预训练模型&对联模型（单方向）&赏析模型\\\midrule
 Transformer层数&6&6&6\\
 隐藏维度／注意力头数&256／4&256／4&256／4\\
 前馈网络中间维度&1024&1024&1024\\
 字符词表规模&13,232&13,232&15,829\\
 最大上下文长度&128&128&1,536\\
 可训练参数量&8,154,624&34,759,428&9,179,904\\
 词汇残差表规模&不使用&202,585词条&不使用\\\bottomrule
 \end{tabularx}
 \par\vspace{4pt}\begin{minipage}{\linewidth}\footnotesize 注：参数量列仅统计可训练参数。对联分支另使用147,337,984参数的冻结ChineseBERT；正反向生成主干分别训练。\end{minipage}
\end{table}

输入标识的形状为$[B,T]$，嵌入后变为$[B,T,256]$；多头注意力的$Q,K,V$分别为$[B,4,T,64]$，注意力分数为$[B,4,T,T]$。多头结果拼接并经过六层变换，最终输出$[B,T,13232]$的logits。图~\ref{fig:base}同时展示主干结构与单层内部计算，明确归一化、注意力、前馈和残差的连接关系。

\begin{figure}[!htbp]
 \centering% Original vector drawing of the implemented Pre-LN model; not a copied paper figure.
\begin{tikzpicture}[x=1cm,y=1cm,
 font=\fontsize{9.2}{11}\selectfont,
 layer/.style={draw=black!70,rounded corners=1pt,line width=.55pt,minimum height=.46cm,align=center,inner sep=3pt},
 tok/.style={draw=black!65,fill=white,minimum width=.58cm,minimum height=.38cm,inner sep=1pt},
 ar/.style={-{Stealth[length=1.8mm]},line width=.6pt},
 plus/.style={circle,draw=black!75,fill=white,inner sep=0pt,minimum size=.38cm}]
% (a) Token-wise decoder stack
\node[anchor=west,font=\small\bfseries] at (0,9.3) {(a) 自回归GPT主干};
\foreach \x/\t in {1/x_1,2/x_2,3/\cdots,4/x_T}{\node[tok] at (\x,.15) {$\t$};}
\node[layer,fill=green!12,minimum width=3.8cm] (emb) at (2.5,1.05) {Token Embedding $E$};
\foreach \x in {1,2,4}{\draw[ar] (\x,.36)--(\x,.81);}
\node[plus] (addp) at (2.5,1.92) {$+$};
\draw[ar] (emb)--(addp);
\node[layer,fill=green!12,minimum width=1.15cm] (pos) at (5.18,1.92) {$P_{1:T}$};
\draw[ar] (pos)--(addp);
% Embedding tensor, shown as a grid rather than a prose process box.
\foreach \i in {0,...,5}{\foreach \j in {0,...,2}{
 \draw[fill=modelblue!35,draw=white,line width=.35pt] (1.51+.33*\i,2.43+.10*\j) rectangle ++(.32,.09);
}}
\node[anchor=west] at (3.75,2.58) {$H^{(0)}\in\mathbb R^{T\times256}$};
\draw[ar] (addp)--(2.5,2.41);
\draw[rounded corners=3pt,fill=modelblue!4,draw=black!60] (.45,3.0) rectangle (4.55,5.6);
\node[layer,fill=modelblue!16,minimum width=3.5cm] at (2.5,3.5) {Pre-LN Block 1};
\node at (2.5,4.28) {$\vdots$};
\node[layer,fill=modelblue!16,minimum width=3.5cm] at (2.5,5.08) {Pre-LN Block 6};
\draw[ar] (2.5,2.75)--(2.5,3.26);
\draw[ar] (2.5,3.75)--(2.5,4.0);\draw[ar] (2.5,4.54)--(2.5,4.84);
\node[anchor=west] at (4.68,4.28) {$6\times$};
\node[layer,fill=black!6,minimum width=3.8cm] (norm) at (2.5,6.12) {LayerNorm};
\node[layer,fill=green!12,minimum width=3.8cm] (head) at (2.5,6.97) {Linear $E^{\mathsf T}$};
\node[layer,fill=orange!12,minimum width=3.8cm] (soft) at (2.5,7.80) {Softmax};
\draw[ar] (2.5,5.33)--(norm);\draw[ar] (norm)--(head);\draw[ar] (head)--(soft);
\draw[ar] (soft)--(2.5,8.35);\node[anchor=south] at (2.5,8.35) {$p(x_{t+1}\mid x_{\le t})$};
\draw[ar,dashed,draw=green!40!black] (emb.west)--(.05,1.05)--(.05,6.97)--(head.west);
\node[rotate=90,fill=white,inner sep=2pt,text=green!35!black] at (.05,6.02) {weight tying};
% (b) Exact Pre-LN residual block, expanded as a neural module.
\node[anchor=west,font=\small\bfseries] at (7.1,9.3) {(b) Pre-LN Block内部结构};
\draw[rounded corners=3pt,draw=black!55,fill=black!2] (7.25,.75) rectangle (12.75,8.15);
\node (hin) at (9.5,.35) {$H^{(\ell-1)}$};
\node[layer,fill=orange!12,minimum width=3.1cm] (ln1) at (9.5,1.45) {LayerNorm};
\node[layer,fill=modelblue!18,minimum width=3.1cm,minimum height=.85cm] (att) at (9.5,2.63) {Masked Multi-Head\\Self-Attention};
\node[plus] (p1) at (9.5,3.72) {$+$};
\node[layer,fill=orange!12,minimum width=3.1cm] (ln2) at (9.5,4.68) {LayerNorm};
\node[layer,fill=purple!11,minimum width=3.1cm,minimum height=.8cm] (ffn) at (9.5,5.83) {Feed-Forward\\$256\to1024\to256$};
\node[plus] (p2) at (9.5,7.0) {$+$};
\node (hout) at (9.5,8.42) {$H^{(\ell)}$};
\draw[ar] (hin)--(ln1);\draw[ar] (ln1)--(att);\draw[ar] (att)--(p1);\draw[ar] (p1)--(ln2);\draw[ar] (ln2)--(ffn);\draw[ar] (ffn)--(p2);\draw[ar] (p2)--(hout);
\draw[ar] (9.5,.92)--(12.1,.92)|-(p1.east);
\draw[ar] (9.5,4.18)--(12.1,4.18)|-(p2.east);
\node[align=center] at (9.5,-.25) {4 heads\quad $d_k=d_v=64$};
\end{tikzpicture}

 \caption{字符级GPT网络架构。(a)词元、嵌入张量、六层解码器与输出分布；虚线表示嵌入和输出投影的权重共享。(b)展开单层Pre-LN计算，圆形加法节点与旁路表示两组残差连接。}
 \label{fig:base}
\end{figure}

基础模型采用长度128的可学习位置表、0.1的Dropout和$10^{-5}$的归一化数值稳定项。注意力既实现显式计算，也与缩放点积注意力接口进行等价性核验；正式训练使用后者提高效率。

\subsection{古典语料构建与数据划分}

预训练文本来自\href{https://github.com/chinese-poetry/chinese-poetry}{chinese-poetry}的唐诗、宋诗和宋词。古文取自\href{https://github.com/NiuTrans/Classical-Modern}{NiuTrans的Classical-Modern语料}。诗词提供韵文句式与常见意象，古文补充叙事、议论和较长语境。两类材料共同服务于古典语言分布学习，不直接提供对联配对监督。

文本处理包括Unicode NFC规范化、\href{https://github.com/BYVoid/OpenCC}{OpenCC}繁简转换、异常记录过滤和重复检查，同时保留标点与换行。古文按书目筛选，采用120部允许书目并排除207部其他书目。划分以作品或书籍为单位组织，避免同一作品的不同窗口跨训练集和验证集；另检查精确及部分近重复，并排查下游保留对联的原文子串。

\begin{table}[htbp]
 \centering\small
 \caption{预训练语料规模与目标采样比例}\label{tab:corpus}
 \begin{tabularx}{\linewidth}{Yrrrr}
 \toprule 来源&训练作品／章节&验证作品／章节&训练正文词元&目标比例\\\midrule
 唐诗&47,922&984&3,005,923&30\%\\
 宋诗&245,840&5,001&15,601,631&30\%\\
 宋词&19,445&372&1,693,281&10\%\\
 古文&5,071&50&10,819,250&30\%\\\midrule
 合计&318,278&6,407&31,120,085&100\%\\\bottomrule
 \end{tabularx}
\end{table}

验证集含630,913个正文词元。表~\ref{tab:corpus}的正文计数包含保留的标点和换行，统计对象为规范化正文。字符词表取自预训练训练集与对联训练集的字符并集，字表构建使用训练侧数据。

由于宋诗规模占比较高，直接遍历会使其主导训练。采用表中的有放回混合采样，使唐诗和古文得到更充分曝光。训练按正文词元预算停止。

\subsection{样本组织与预训练策略}

每篇文本按BOS、正文和EOS组织，在篇内构造长度128、步长128的窗口并保留尾部，不跨篇拼接。输入与标签右移一位：例如输入“BOS、春、风、又、绿”，对应目标为“春、风、又、绿、EOS”。较短窗口右侧补PAD，其标签不参与损失。作品级划分在切窗前完成，同一作品的所有窗口继承其集合归属。

微批次大小为32，累计至少8192个有效目标后执行一次优化器更新。有效目标数包含正文和EOS，因此不同于正文采样预算。训练以实际累计目标量归一化梯度，并将范数裁剪至1。AdamW的$\beta$为$(0.9,0.95)$，主要线性变换矩阵权重衰减为0.1，嵌入和归一化等按参数分组免衰减。

学习率在前100万正文词元内预热至$3\times10^{-4}$，随后余弦下降至$3\times10^{-5}$。本次训练处理30,004,916个正文词元、30,329,796个有效监督目标，完成3,337次更新。唐诗、宋诗、宋词和古文的实际正文占比分别为30.015\%、30.103\%、9.925\%和29.957\%，与目标比例接近。预训练运行记录耗时约366.48秒，峰值已分配GPU内存约914.95 MiB，统计对象为预训练进程。

\subsection{验证结果与生成行为}

完整验证按既定来源比例加权的交叉熵从9.5545降至4.7773，其指数从14,107.76降至118.79。这里先加权各来源交叉熵再取指数，并非将所有验证词元直接合并计算的困惑度。图~\ref{fig:loss}保留固定训练探针和完整验证的原始评估点，不进行曲线平滑。

\begin{figure}[htbp]
 \centering\begin{tikzpicture}
\begin{axis}[width=.94\linewidth,height=6.4cm,xmin=0,xmax=30.7,ymin=4,ymax=10,
 xlabel={累计训练正文词元数（百万）},ylabel={平均交叉熵},
 tick label style={font=\small},label style={font=\small},
 grid=major,grid style={black!10},axis line style={black!60},
 legend style={font=\small,at={(.98,.98)},anchor=north east,draw=none,fill=white},
 xtick={0,5,10,15,20,25,30},ytick={4,5,6,7,8,9,10}]
\addplot[color=modelblue,thick,mark=triangle*,mark size=2.2pt] table[x=million,y=ce,col sep=comma]{figures/full_validation.csv};
\addlegendentry{完整验证集}
\addplot[color=black!65,dashed,mark=o,mark size=1.4pt] table[x=million,y=ce,col sep=comma]{figures/train_probe.csv};
\addlegendentry{固定训练探针}
\end{axis}
\end{tikzpicture}

 \caption{预训练交叉熵变化。完整验证按约500万正文词元间隔评估；固定训练探针在关闭Dropout后计算。两条曲线分别连接各固定集合的原始评估点。}
 \label{fig:loss}
\end{figure}

\begin{table}[htbp]
 \centering\small
 \caption{预训练结束时的完整验证结果}\label{tab:pretrainmetrics}
 \begin{tabular}{lrr}
 \toprule 语料来源&交叉熵&困惑度\\\midrule
 唐诗&4.5193&91.77\\宋诗&4.7230&112.50\\宋词&4.6349&103.01\\古文&5.1372&170.24\\
 来源加权结果&4.7773&118.79\\\bottomrule
 \end{tabular}
\end{table}

\Needspace{5\baselineskip}
表~\ref{tab:pretrainmetrics}进一步给出各来源的验证结果。在固定提示“江上”之后，预训练模型生成“春风满，江边柳色新。不知春色好，不是故人稀”，表现出诗词常见意象和局部句式。另一些输入出现“今朝不可知”“今朝不可见”等反复。这些输出呈现了模型对古典意象和句式的学习，也记录了重复生成现象。

\subsection{预训练对任务迁移的作用}

为检验预训练的作用，保持8.15M参数GPT、400,290条对联训练配对及监督预算一致，比较随机初始化和预训练初始化。三个成对种子均训练3轮，每个分支呈现1,200,870个样本、监督9,192,768个下联及EOS词元，以验证交叉熵选点。生成在2,233个去重测试上联上采用贪心解码。表~\ref{tab:transfer}给出三种子均值与样本标准差。

\begin{table}[htbp]
\centering\small
\caption{相同监督预算下的预训练迁移对照}\label{tab:transfer}
\begin{tabularx}{\linewidth}{Yrr}
\toprule 指标&随机初始化&预训练后微调\\\midrule
下联及EOS交叉熵&$3.7478\pm0.0047$&$3.5237\pm0.0031$\\
字符预测准确率（\%）&$22.9435\pm0.0699$&$25.6275\pm0.0838$\\
生成chrF&$2.5320\pm0.0598$&$2.7212\pm0.0167$\\\bottomrule
\end{tabularx}
\end{table}

预训练组的目标损失更低、字符准确率和参考重合更高，三个微调种子的变化方向一致。预训练由此为配对任务提供了有利的参数起点；下一节以原框架直接微调模型为基线，进一步研究其词汇和形式对应问题。

\subsection{实现正确性核验}

实验核验了未来词元扰动下的因果性、右侧补齐不变性、显式注意力与高效接口的一致性，以及小样本过拟合、检查点恢复和权重重载。上述检查分别核对信息时序、监督对齐与状态恢复。通过这些检查后，采用训练所得的共享主干开展后续任务微调。
